\documentclass[12pt,a4paper]{article}

\usepackage[utf8]{inputenc}
\usepackage[english]{babel}
\usepackage{graphicx}
\usepackage{float}
\usepackage{geometry}
\usepackage{amsmath}

\geometry{top=2.5cm, bottom=2.5cm, left=2.5cm, right=2.5cm}

\title{Main Results: Magnetic Coupling in 7-AGNRs}
\author{}
\date{\today}

\begin{document}

\maketitle

\section{Introduction}
This document summarizes the main results obtained for the calculation of the magnetic exchange coupling $J$ in finite graphene nanoribbons (7-AGNRs) of different lengths, comparing the Tight-Binding model (effective Hubbard dimer) with the Mean-Field Hubbard (MFH) model.

\section{Geometry}
First, we verify the geometric construction of the system by removing atoms with coordination 1 at the zigzag edges to ensure a properly terminated finite molecule.

\begin{figure}[H]
    \centering
    \includegraphics[width=0.8\textwidth]{test_geom_L4.png}
    \caption{Geometry test for a short 7-AGNR ribbon ($L=4$).}
    \label{fig:geom}
\end{figure}

\section{Spin Polarization (AFM and FM)}
The ground state of the edges exhibits strong localization and spin polarization. Figure \ref{fig:spin} shows the spin polarization profiles calculated using Mean-Field Hubbard (MFH) for the Antiferromagnetic (AFM, open-shell singlet) and Ferromagnetic (FM, triplet) configurations.

\begin{figure}[H]
    \centering
    \includegraphics[width=0.8\textwidth]{spin_polarization_L6.png}
    \caption{Spin distribution for the AFM and FM states calculated in a ribbon of length $L=6$.}
    \label{fig:spin}
\end{figure}

\section{Magnetic Coupling $J$ Comparison}
We calculate the magnetic coupling $J$ using two methods:
\begin{itemize}
    \item The effective Hubbard dimer rule ($J = 4t_{\mathrm{eff}}^2 / U_{\mathrm{eff}}$) derived from Tight-Binding.
    \item The energy difference between the FM and AFM states obtained from the MFH self-consistent cycle ($J = E_{\mathrm{FM}} - E_{\mathrm{AFM}}$).
\end{itemize}

The plot also includes experimental references, such as the superconducting gap $2\Delta$ for Nb(110) and the typical thermal energy scale of an STM operating at 1.2 K.

\begin{figure}[H]
    \centering
    \includegraphics[width=0.9\textwidth]{J_comparison_vs_L.png}
    \caption{Comparison of the exchange coupling $J$ extracted by both models as a function of the length $L$ (number of DBBA monomers).}
    \label{fig:comparison}
\end{figure}

\section{Energy Difference Between Models}
To quantitatively observe the deviation between both methods, we analyze the absolute difference of the calculated couplings.

\begin{figure}[H]
    \centering
    \includegraphics[width=0.9\textwidth]{J_diff_vs_L.png}
    \caption{Absolute difference $|\Delta J| = |J_{\mathrm{MFH}} - J_{\mathrm{Dimer}}|$ as a function of length $L$.}
    \label{fig:diff}
\end{figure}

\end{document}
